% Part: history
% Chapter: biographies
% Section: bertrand-russell
%READERNOTE{SOL6-A275: मूळातील सहा महिन्यांचा उल्लेख शिक्षेचा आहे; प्रत्यक्ष सुटका १४ सप्टेंबर १९१८ रोजी झाली. तसेच आयुष्यभर बिनशर्त शांततावाद हा दावा मागे घेतला: मॅकमास्टरच्या 30909 या पत्रनोंदीतील रसेल यांच्या स्वतःच्या विधानात काही परिस्थितींमध्ये बलप्रयोग समर्थनीय मानला आहे. स्रोत: https://russell-letters.mcmaster.ca/chronology आणि https://bracers.mcmaster.ca/30909.}

\documentclass[../../../include/open-logic-section]{subfiles}

\begin{document}

\olfileid{his}{bio}{rus} 

\olsection{बर्ट्रंड रसेल}

\olphoto{russell-bertrand}{बर्ट्रंड रसेल}

बर्ट्रंड रसेल यांना आधुनिक विश्लेषणात्मक तत्त्वज्ञानाच्या संस्थापकांपैकी एक मानले जाते. १८ मे १८७२ रोजी जन्मलेले रसेल तत्त्वज्ञान आणि तर्कशास्त्रातील कामासाठी तर ओळखले जातच; याशिवाय त्यांनी अनेक विषयांवर सर्वसामान्य वाचकांसाठी पुस्तके लिहिली. आयुष्यभर ते सक्रिय राजकीय कार्यकर्तेही होते.

रसेल यांचा जन्म वेल्समधील मॉनमथशरच्या ट्रेलेच येथे झाला. त्यांचे आईवडील ब्रिटिश उमराव वर्गातील होते. ते स्वतंत्र विचारांचे होते आणि त्या काळी बोस्टनमधील क्रांतिकारी विचारांच्या लोकांशीही त्यांची मैत्री होती. दुर्दैवाने रसेल लहान असतानाच त्यांच्या आईवडिलांचे निधन झाले; म्हणून ते आजीआजोबांकडे राहायला गेले. तेथे त्यांचे धार्मिक संस्कारांत संगोपन झाले, जे त्यांच्या आईवडिलांना कोणत्याही परिस्थितीत टाळायचे होते. नैतिकतेच्या सर्व बाबतीत त्यांची आजी अतिशय कडक होती. किशोरवयात त्यांचे बहुतांश शिक्षण खाजगी शिक्षकांकडून घरीच झाले.

विश्लेषणात्मक तत्त्वज्ञानावर, विशेषतः तर्कशास्त्रावर, रसेल यांचा प्रभाव फार मोठा आहे. त्यांनी केंब्रिजच्या ट्रिनिटी कॉलेजमध्ये गणित आणि तत्त्वज्ञान शिकले; तेथे गणितज्ञ आणि तत्त्वज्ञ ॲल्फ्रेड नॉर्थ व्हाइटहेड यांचा त्यांच्यावर प्रभाव पडला. १९१० मध्ये रसेल आणि व्हाइटहेड यांनी \emph{Principia Mathematica} चा पहिला खंड प्रकाशित केला. त्यात गणिताचे तर्कशास्त्रात न्यूनीकरण करता येते, या भूमिकेचे त्यांनी समर्थन केले. पुढे रसेल यांनी शेकडो पुस्तके, निबंध आणि राजकीय पत्रके प्रकाशित केली. १९५० मध्ये त्यांना साहित्यातील नोबेल पारितोषिक मिळाले.

रसेल राजकारणात आणि सामाजिक चळवळींत खोलवर सहभागी होते. पहिल्या महायुद्धात शांततावादी कृती आणि निषेध यांमुळे त्यांना अटक झाली आणि सहा महिन्यांच्या तुरुंगवासाची शिक्षा झाली. तुरुंगात त्यांना वाचता आणि लिहिता येत असे; हा अनुभव “बराच सुखावह” वाटल्याचा त्यांचा दावा होता. पुढेही त्यांनी शांततेसाठी काम केले, मात्र सर्व परिस्थितींमध्ये युद्धाला विरोध करावा अशी त्यांची भूमिका नव्हती. १९६१ मध्ये अण्वस्त्रनिःशस्त्रीकरणाच्या सभेत सहभागी झाल्यामुळे त्यांना पुन्हा तुरुंगवास झाला. १९४८ मध्ये ते विमान अपघातातूनही बचावले. या स्रोताच्या कथनानुसार, तेव्हा धूम्रपानासाठी राखीव विभागात बसलेले लोकच बचावले; म्हणून धूम्रपानामुळेच आपला जीव वाचला, असा रसेल यांचा दावा होता. रसेल यांनी चार विवाह केले, पण विवाहबाह्य संबंध ठेवत असल्याची त्यांची ख्याती होती. २ फेब्रुवारी १९७० रोजी वयाच्या ९७व्या वर्षी वेल्समधील पेनरिनड्यूड्रेथ येथे त्यांचे निधन झाले.

\begin{reading}
रसेल यांनी १८७२ ते १९६७ या कालखंडातील आपले जीवन सांगणारे तीन भागांचे आत्मचरित्र लिहिले \citep{Russell1967,Russell1968,Russell1969}. मॅकमास्टर विद्यापीठातील बर्ट्रंड रसेल संशोधन केंद्रात त्यांच्या संग्रहित दस्तऐवजांचा ठेवा आहे. त्यांच्या संकलित लेखनाचे खंड, शोधता येणाऱ्या अनुक्रमणिका आणि संग्रहण प्रकल्पांविषयी केंद्राच्या संकेतस्थळावर \citet{Duncan2015} माहिती आहे. रसेल यांचा \emph{On
  Denoting} \citep{Russell1905} हा लेख विसाव्या शतकातील विश्लेषणात्मक तत्त्वज्ञानातील अभिजात लेख मानला जातो.

रसेल यांच्यावरील स्टॅनफर्ड एन्सायक्लोपीडिया ऑफ फिलॉसॉफीमधील नोंदीत \citep{Irvine2015} इच्छा आणि राजकीय सिद्धांताविषयी रसेल बोलत असल्याच्या ध्वनिमुद्रिका आहेत. रसेल यांच्या अनेक चित्रित मुलाखती ऑनलाइन उपलब्ध आहेत. उदाहरणार्थ, धूम्रपान आणि विमान अपघातातील सहभागाविषयी त्यांना बोलताना पाहण्यासाठी \citet{RussellND} पाहा. त्यांच्या काही पुस्तकांची, त्यांत \emph{Introduction to Mathematical Philosophy} चा समावेश आहे, विनामूल्य ध्वनिपुस्तके \citet{LibriVoxND} येथे उपलब्ध आहेत.
\end{reading}

\end{document}
